% ПРИЛОЖЕНИЯ — без номер.
% По А.3: изходен текст на програми, формати на структури от данни, екранни форми,
% получени резултати.
\section*{ПРИЛОЖЕНИЯ}
\label{sec:appendix}
\addcontentsline{toc}{section}{ПРИЛОЖЕНИЯ}

\subsection*{Приложение А. Формат на записа с товарните таблици}
\addcontentsline{toc}{subsection}{Приложение А. Формат на записа с товарните таблици}

Наборът е един обект с три полета: \code{meta} с описанието на мерните системи,
\code{walls} с конфигурациите на катерачните стени, ключирани по височина и брой нива,
и \code{boulder} с редовете за боулдер стени. Листинг~\ref{lst:datafmt} показва един
съкратен ред; пълният набор съдържа 26\,600 числови стойности.

\begin{lstlisting}[language=JSONlang,caption={Съкратен запис на една конфигурация},label={lst:datafmt}]
"W_11_2": {
  "height": 11, "levels": 2,
  "lhEU": { "1": 5.5, "2": 10.5 },
  "lhUS": { "1": 18,  "2": 34.4 },
  "rows": [
    { "x": 1, "a": 4, "lvl": 1,
      "eu": { "RZ0DL": -5.38, "RX0DL": -0.69, "RX1DL": 0.29, "RX2DL": 0.39,
              "RZ0LL": -8,    "RX0LL": -1.24, "RX1LL": 2.26, "RX2LL": 0.54,
              "LX1DL": 1.16,  "LX2DL": 1.56,  "LX1LL": 3.955, "LX2LL": 0.945 },
      "us": { "RZ0DL": -1210.5, "RX0DL": -155.25, "...": "..." } }
  ]
}
\end{lstlisting}

Полетата \code{lhEU} и \code{lhUS} дават височините на нивата на закрепване и се
използват както за надписите в таблицата с резултатите, така и за построяването на
схемата. Полетата \code{a} и \code{x} са разстоянието между колоните и наклонът,
\code{lvl} е номерът на нивото, за което важи редът.
\newpage
\subsection*{Приложение Б. Формат на записа за проект}
\addcontentsline{toc}{subsection}{Приложение Б. Формат на записа за проект}

\begin{lstlisting}[language=JSONlang,caption={Публично представяне на един проект},label={lst:projfmt}]
{
  "id": "6a807fd62bbc9f4b753e1834",
  "name": "Sofia Gym - main hall",
  "tags": ["gym", "EU", "zone-1"],
  "properties": [ { "key": "Client", "value": "Acme Climbing Ltd" },
                  { "key": "Ref",    "value": "WT-2026-001" } ],
  "input": { "units": "EU", "type": "wall", "height": 12, "levels": 3,
             "span": 6, "overhang": 1, "force": 1, "factored": false,
             "baseDetail": "CF-02", "attachmentDetail": "CW-03",
             "loadsRequestedBySolution": { "single": true, "beams": false } },
  "snapshot": { "title": "Climbing wall 12 m", "unit": "kN",
                "governing": 5.15, "verdict": "ok" },
  "createdAt": "2026-08-15T15:11:02.317Z",
  "updatedAt": "2026-08-15T15:11:02.317Z"
}
\end{lstlisting}

Полетата \code{baseDetail} и \code{attachmentDetail} са идентификаторите на избрания
типов детайл за подовата плоча и за метода на закрепване (§\ref{subsec:attachimpl}).
Полето \code{loadsRequestedBySolution} пази кое от двете схеми на закрепване
(§\ref{subsec:ug-options}) потребителят е разгърнал, за да бъде възпроизведено същото
разгънато състояние при отваряне на записания проект.
\newpage
\subsection*{Приложение В. Признак за нарушение при проверката на разположението}
\addcontentsline{toc}{subsection}{Приложение В. Признак за нарушение при проверката на разположението}

Листинг~\ref{lst:probe} съдържа ядрото на проверката по §\ref{subsec:howcheck}:
обхождането на документа и трите изключения, обосновани там.

\begin{lstlisting}[caption={Откриване на елементите, които разширяват страницата},label={lst:probe}]
function scan(doc, win) {
  const vw = doc.documentElement.clientWidth;
  const inScroller = (el) => {
    let p = el.parentElement;
    while (p && p !== doc.body) {
      const o = win.getComputedStyle(p).overflowX;
      if (["auto", "scroll", "hidden", "clip"].includes(o)) return true;
      p = p.parentElement;
    }
    return false;
  };
  const bad = [];
  for (const el of doc.querySelectorAll("body *")) {
    const r  = el.getBoundingClientRect();
    if (!r.width && !r.height) continue;
    const cs = win.getComputedStyle(el);
    if (cs.position === "fixed" || cs.display === "none") continue;
    if ((r.right > vw + 1 || r.left < -1) && !inScroller(el)) bad.push(el);
  }
  return { violations: bad,
           documentOverflow: doc.documentElement.scrollWidth - vw };
}
\end{lstlisting}

\subsection*{Приложение Г. Отстраняване на основния дефект}
\addcontentsline{toc}{subsection}{Приложение Г. Отстраняване на основния дефект}

Клас Д1 от \tabref{tab:taxonomy}, първопричината на всичките 17 наблюдавани
нарушения, се отстранява с едно правило:

\begin{lstlisting}[language=CSS,caption={Основният блок не приема вътрешната си минимална ширина},label={lst:fix}]
/* `body` is a column flex container and every page's <main> is centred with
   `margin-inline: auto`. An auto cross-axis margin cancels `align-items:
   stretch`, so <main> falls back to fit-content sizing - which is at least
   its min-content width. Pinning the width to 100% keeps <main> exactly
   viewport-wide and leaves overflow to the inner scroll containers. */
body > main { flex: 1 0 auto; width: 100%; min-width: 0; }
\end{lstlisting}

\subsection*{Приложение Д. Отчет от изпитанието на приложния програмен интерфейс}
\addcontentsline{toc}{subsection}{Приложение Д. Отчет от изпитанието}

Екранните форми са дадени в Раздел~\ref{sec:userguide} на мястото, където съответната
функция се описва, и не се повтарят тук.

\begin{lstlisting}[language={},caption={Изход от изпълнението на изпитанието (съкратен)},label={lst:apilog}]
  ok - health ok                        ok - login succeeds
  ok - register rejects bad email       ok - login returns a bearer token
  ok - register rejects short password  ok - bearer token authenticates
  ok - register creates user            ok - invalid bearer token -> 401
  ok - register sets auth cookie        ok - create / read / update project
  ok - me returns current user          ok - attachment selections stored
  ok - me is 401 without cookie         ok - filter by tag, sort by name
  ok - duplicate email -> 409           ok - other user's project -> 404
  ok - wrong password -> 401            ok - other user's support -> 404
                                         ok - delete project
                                         ok - deleted project gone
                                         ok - logout

ALL 36 CHECKS PASSED
\end{lstlisting}

\subsection*{Приложение Е. Указания за изграждане и стартиране}
\addcontentsline{toc}{subsection}{Приложение Е. Изграждане и стартиране}

\begin{lstlisting}[language=bash,caption={Стартиране на услугата и изпитанията},label={lst:run}]
# Услуга (изисква server/.env с MONGODB_URI и JWT_SECRET)
cd server && npm install && npm start          # http://localhost:8787

# Услуга с база в паметта, за разработка и изпитания
cd server && node test/serve-fake.js

# Изпитание на приложния програмен интерфейс
cd server && node test/api.test.js
\end{lstlisting}

Уеб клиентът не изисква изграждане: услугата обслужва директорията \code{website/}
като статични файлове, така че при промяна е достатъчно презареждане на страницата.
